\chapterbackmatter{Solutions to Supplementary Exercises}

\begin{enumerate}
\SupplementarySolution{1}{One Phase-Space Time Slice}
Let \(\delta t>0\), \(q'=q(t+\delta t)\), and \(q=q(t)\).  To the
accuracy needed in a time-sliced product,
\[
 e^{-iH\delta t}
 =e^{-iP^2\delta t/(2m)}e^{-iV(Q)\delta t}
 +O(\delta t^2).
\]
Inserting \(\mathbf1=\int dp\,\ket p\bra p\), with
\(\braket{q}{p}=(2\pi)^{-1/2}e^{ipq}\), gives
\[
 \braket{q'}{e^{-iH\delta t}q}
 =\int\frac{dp}{2\pi}
 \exp\left\{ip(q'-q)-i\delta t
 \left[\frac{p^2}{2m}+V(q)\right]\right\}+O(\delta t^2).
\]
The potential is evaluated at the right endpoint of the ket because of the
displayed operator ordering.  A midpoint prescription instead represents
Weyl ordering; the distinction matters for nontrivial products of \(P\)
and \(Q\), but vanishes here in the continuum limit.

\SupplementarySolution{2}{Integrating Out a Quadratic Momentum}
For each slice, complete the square:
\[
 p\frac{\Delta q}{\delta t}-\frac{p^2}{2m}
 =-\frac1{2m}\left(p-m\frac{\Delta q}{\delta t}\right)^2
 +\frac m2\left(\frac{\Delta q}{\delta t}\right)^2.
\]
Thus
\[
 \int\frac{dp}{2\pi}e^{i\delta t[p\Delta q/\delta t-p^2/(2m)]}
 =\left(\frac{m}{2\pi i\delta t}\right)^{1/2}
 e^{im(\Delta q)^2/(2\delta t)}.
\]
Multiplying the slices and taking their continuum limit yields
\[
 K(q_f,t_f;q_i,t_i)
 =\int_{q_i}^{q_f}\mathcal Dq\,
 \exp\left\{i\int_{t_i}^{t_f}dt
 \left[\frac m2\dot q^2-V(q)\right]\right\}.
\]
The coefficient of \(p^2\) is the constant \(1/m\), so the product of
Gaussian determinants depends only on the slicing and total time.  A
coordinate-dependent kinetic coefficient would instead leave a
field-dependent measure factor.

\SupplementarySolution{3}{The Schrödinger Equation from One Step}
Write \(y=x-\xi\) in
\[
 \psi(x,t+\delta t)=
 \left(\frac{m}{2\pi i\delta t}\right)^{1/2}
 \int d\xi\,
 e^{im\xi^2/(2\delta t)}
 e^{-i\delta t V(x-\xi)}\psi(x-\xi,t).
\]
The oscillatory Gaussian localizes \(\xi=O(\sqrt{\delta t})\).  Expanding
the last two factors and using
\[
 \langle\xi\rangle=0,\qquad
 \langle\xi^2\rangle=\frac{i\delta t}{m},
\]
gives
\[
 \psi(x,t+\delta t)=\psi
 +\frac{i\delta t}{2m}\partial_x^2\psi
 -i\delta t\,V(x)\psi+O(\delta t^{3/2}).
\]
After subtracting \(\psi(x,t)\), dividing by \(\delta t\), and multiplying
by \(i\), one obtains
\[
 i\partial_t\psi=
 \left[-\frac1{2m}\partial_x^2+V(x)\right]\psi.
\]
The normalization of the short-time kernel is essential: it makes the
zeroth Gaussian moment equal to one.

\SupplementarySolution{4}{The Free-Particle Kernel}
Successive convolution of the \(N\) identical Gaussian slices gives
\[
 K_0(x_f,T;x_i,0)=
 \left(\frac{m}{2\pi iT}\right)^{1/2}
 \exp\left[\frac{im(x_f-x_i)^2}{2T}\right],
\qquad T>0.
\]
Equivalently, this is the momentum integral
\[
 \int\frac{dp}{2\pi}
 e^{ip(x_f-x_i)-ip^2T/(2m)}.
\]
As \(T\to0^+\), the latter representation converges distributionally to
\(\delta(x_f-x_i)\), fixing the branch of the square root.  In units with
\(\hbar=1\), \(m/T\) has dimension length\(^{-2}\), so the prefactor has
dimension length\(^{-1}\), as required for a one-dimensional position
kernel.

\SupplementarySolution{5}{Composition Fixes the Normalization}
Write
\[
 K_0(x_2,T;x_0,0)=A(T)
 e^{im(x_2-x_0)^2/(2T)}.
\]
Composition at \(T=T_1+T_2\) requires
\[
 A(T)=A(T_2)A(T_1)
 \left(\frac{2\pi iT_1T_2}{mT}\right)^{1/2}.
\]
Hence \(B(T)\equiv A(T)\sqrt T\) is constant, and substitution back into
the same equation gives \(B=\sqrt{m/(2\pi i)}\).  Therefore
\[
 A(T)=\left(\frac{m}{2\pi iT}\right)^{1/2}.
\]
The sign and phase of the root are chosen continuously from \(T=0^+\);
with this choice \(K_0\) approaches the delta distribution rather than its
negative.

\SupplementarySolution{6}{The Dirichlet Fluctuation Spectrum}
Set \(x=x_{\rm cl}+\eta\), where
\(\eta(0)=\eta(T)=0\).  The fluctuation action is
\[
 S_{\rm fl}=\frac m2\int_0^Tdt\,\dot\eta^2
 =\frac12\int_0^Tdt\,\eta(-m\partial_t^2)\eta.
\]
The normalized eigenfunctions and eigenvalues are
\[
 u_n(t)=\sqrt{\frac2T}\sin\frac{n\pi t}{T},
\qquad
 \lambda_n=m\left(\frac{n\pi}{T}\right)^2,
\qquad n=1,2,\ldots.
\]
The Gaussian fluctuation integral is formally
\((\det[-m\partial_t^2])^{-1/2}\).  Comparing determinant ratios at two
times, or using zeta regularization with \(\zeta(0)=-1/2\), leaves the
scaling \(T^{-1/2}\).  The absolute constant is not fixed by the continuum
determinant alone; the short-time delta normalization supplies
\(\sqrt{m/(2\pi i)}\).

\SupplementarySolution{7}{The Oscillator Feynman Green Function}
Define
\[
 G_F(\tau)=\int\frac{dE}{2\pi}
 \frac{e^{-iE\tau}}{E^2-\omega^2+i0}.
\]
For \(\tau>0\) the contour closes below and selects the pole near
\(+\omega\); for \(\tau<0\) it closes above and selects the pole near
\(-\omega\).  The result is
\[
 G_F(\tau)=\frac{i}{2\omega}e^{-i\omega|\tau|}.
\]
Away from \(\tau=0\), it obeys the homogeneous oscillator equation.  Its
first derivative has jump
\(G_F'(0^+)-G_F'(0^-)=1\), and consequently
\[
 (\partial_\tau^2+\omega^2)G_F(\tau)=\delta(\tau).
\]
The positive-frequency mode propagates forward and the negative-frequency
mode backward, which is exactly oscillator time ordering.

\SupplementarySolution{8}{Vacuum Survival under an External Force}
The forced action remains Gaussian.  Completing its functional square gives
\[
 \frac{Z[f]}{Z[0]}
 =\exp\left[-\frac{i}{2}\int dt\,dt'\,
 f(t)G_F(t-t')f(t')\right],
\]
with the convention for \(G_F\) in S.9.7.  Only the delta-function part of
the pole prescription contributes to the modulus.  For
\(\widetilde f(E)=\int dt\,e^{iEt}f(t)\) and real \(f\),
\[
 \boxed{\quad
 P_{0\to0}=\left|\frac{Z[f]}{Z[0]}\right|^2
 =\exp\left[-\frac{|\widetilde f(\omega)|^2}{2\omega}\right].
 \quad}
\]
The exponent is the mean occupation number of the coherent final state.
It is dimensionless, nonnegative, and vanishes for an adiabatic force with
no Fourier support at the oscillator frequency.

\SupplementarySolution{9}{The Free Scalar Source Functional}
After integrating by parts, the free action is quadratic.  In momentum
space it reads
\[
 S_0+S_J=
 -\frac12\int\frac{d^4p}{(2\pi)^4}
 \widetilde\phi(-p)(p^2+m^2-i0)\widetilde\phi(p)
 +\int\frac{d^4p}{(2\pi)^4}\widetilde J(-p)\widetilde\phi(p).
\]
Shift
\(\widetilde\phi(p)\mapsto\widetilde\phi(p)
-(p^2+m^2-i0)^{-1}\widetilde J(p)\).  Translation invariance of the
measure then yields
\[
 \boxed{\quad
 \frac{Z_0[J]}{Z_0[0]}
 =\exp\left[\frac i2\int\frac{d^4p}{(2\pi)^4}
 \frac{\widetilde J(-p)\widetilde J(p)}
 {p^2+m^2-i0}\right].
 \quad}
\]
The positive infinitesimal is fixed by the vacuum boundary functionals,
not by the algebraic completion of the square.

\SupplementarySolution{10}{Wick Pairings by Differentiation}
Let
\[
 G_F(x-y)\equiv
 \left.\frac1{i^2Z_0[0]}
 \frac{\delta^2Z_0[J]}{\delta J(x)\delta J(y)}
 \right|_{J=0}.
\]
Because the exponent of \(Z_0[J]\) is quadratic, four derivatives can be
nonzero at \(J=0\) only if they differentiate it in pairs.  Therefore
\[
 \begin{split}
 \langle T\phi_1\phi_2\phi_3\phi_4\rangle_0={}&
 G_{F,12}G_{F,34}+G_{F,13}G_{F,24}
 +G_{F,14}G_{F,23}.
 \end{split}
\]
For \(2n\) fields, every complete pairing occurs once.  The number of
terms is
\[
 (2n-1)!!=\frac{(2n)!}{2^n n!}.
\]
An odd number of derivatives leaves one source unpaired and hence gives
zero at \(J=0\).

\SupplementarySolution{11}{The Contact Term in the Green Equation}
Write
\[
 G_F(x-y)=\theta(x^0-y^0)
 \bra{\Omega}\phi(x)\phi(y)\ket{\Omega}
 +\theta(y^0-x^0)
 \bra{\Omega}\phi(y)\phi(x)\ket{\Omega}.
\]
The Klein--Gordon operator annihilates the fields away from equal time.
The first time derivative produces no delta term because
\([\phi(t,\mathbf x),\phi(t,\mathbf y)]=0\).  The second produces
\[
 \delta(x^0-y^0)
 [\dot\phi(x),\phi(y)]_{x^0=y^0}
 =-i\delta(x^0-y^0)\delta^3(\mathbf x-\mathbf y).
\]
With the sign from \(\Box=\nabla^2-\partial_t^2\), the two negatives
cancel, and
\[
 \boxed{\quad
 (\Box_x-m^2)G_F(x-y)=i\delta^4(x-y).
 \quad}
\]
This contact term is the path-integral Green equation in operator form.

\SupplementarySolution{12}{The Scalar Vacuum Wave Functional}
For each \(\mathbf p\), the scalar Hamiltonian is an oscillator of frequency
\(E_{\mathbf p}=\sqrt{\mathbf p^2+m^2}\).  The product of their ground
states is
\[
 \Psi_0[\phi]=\mathcal N
 \exp\left[-\frac12\int d^3p\,
 E_{\mathbf p}\widetilde\phi(\mathbf p)
 \widetilde\phi(-\mathbf p)\right].
\]
In position space the kernel is
\[
 \mathcal E(\mathbf x-\mathbf y)=
 \int\frac{d^3p}{(2\pi)^3}
 E_{\mathbf p}e^{i\mathbf p\cdot(\mathbf x-\mathbf y)},
\]
so it is nonlocal even though the Hamiltonian density is local.  Gaussian
integration gives the equal-time covariance
\[
 \langle\widetilde\phi(\mathbf p)
 \widetilde\phi(\mathbf q)\rangle_0
 =\frac{\delta^3(\mathbf p+\mathbf q)}{2E_{\mathbf p}}.
\]
This inverse-frequency factor agrees with the canonical mode expansion.

\SupplementarySolution{13}{Vacuum Boundaries and the Normalized Ratio}
On a finite interval the path integral contains the boundary factors
\(\Psi_0[\phi_i]\Psi_0[\phi_f]^*\).  Thus
\[
 \langle T\phi(x_1)\cdots\phi(x_n)\rangle
 =\lim_{T\to\infty}
 \frac{\int_{\!-T}^{T}\mathcal D\phi\,
 \Psi_0[\phi_f]^*\phi(x_1)\cdots\phi(x_n)
 e^{iS[\phi]}\Psi_0[\phi_i]}
 {\int_{\!-T}^{T}\mathcal D\phi\,
 \Psi_0[\phi_f]^*e^{iS[\phi]}\Psi_0[\phi_i]}.
\]
Equivalently, adiabatically replace \(H\) by \((1-i\epsilon)H\).  Excited
components at the two ends are exponentially damped, leaving the vacuum.
The denominator removes the vacuum-to-vacuum amplitude, including its
phase and disconnected vacuum bubbles, so the normalized expectation of
\(\mathbf1\) is exactly one.

\SupplementarySolution{14}{No-Production Probability from a Source}
The normalized vacuum amplitude is
\[
 \OutBra{\Omega}\InKet{\Omega;J}
 =\exp\left[\frac i2\int\frac{d^4p}{(2\pi)^4}
 \frac{\widetilde J(-p)\widetilde J(p)}
 {p^2+m^2-i0}\right].
\]
Using
\((u-i0)^{-1}=\operatorname{PV}(1/u)+i\pi\delta(u)\), the principal-value
part is a phase.  Hence
\[
 \boxed{\quad
 P_0=
 \left|\OutBra{\Omega}\InKet{\Omega;J}\right|^2
 =\exp\left[-\int\frac{d^3p}{(2\pi)^3\,2E_{\mathbf p}}
 |\widetilde J(E_{\mathbf p},\mathbf p)|^2\right].
 \quad}
\]
The exponent is minus the mean number of produced quanta.  Only the
on-shell source appears, so a sufficiently slow source produces no
particles.

\SupplementarySolution{15}{Connected Two-Point Graphs from the Path Integral}
At first order,
\[
 \langle T\phi(x)\phi(y)\rangle
 =\frac{\langle T\phi(x)\phi(y)
 [1-i\lambda\int d^4z\,\phi(z)^4/4!]\rangle_0}
 {\langle T[1-i\lambda\int d^4z\,\phi(z)^4/4!]\rangle_0}.
\]
For the connected tadpole, choose the vertex field joined to \(x\) in
\(4\) ways and that joined to \(y\) in \(3\) ways; the remaining two pair.
Thus \(12/4!=1/2\), and
\[
 G_F(x-y)-\frac{i\lambda}{2}\int d^4z\,
 G_F(x-z)G_F(0)G_F(z-y)+O(\lambda^2).
\]
The other numerator contraction pairs \(x\) with \(y\) and contracts the
four vertex fields among themselves in three ways.  It is exactly the
free two-point function times the first-order denominator and therefore
cancels.

\SupplementarySolution{16}{The Tree Four-Point Function}
The first connected contribution is
\[
 -\frac{i\lambda}{4!}\int d^4z\,
 \langle T\phi(x_1)\phi(x_2)\phi(x_3)\phi(x_4)\phi(z)^4\rangle_0.
\]
Connectivity at tree level requires each external field to contract with
a different field at \(z\).  There are \(4!\) bijections, canceling the
factor in the interaction.  Consequently
\[
 \boxed{\quad
 G^{(4)}_{\rm conn}(x_1,\ldots,x_4)
 =-i\lambda\int d^4z\,
 \prod_{a=1}^{4}G_F(x_a-z)+O(\lambda^2).
 \quad}
\]
Contractions that pair external fields among themselves are disconnected;
pure vacuum factors cancel against \(Z[0]\).

\SupplementarySolution{17}{The Vacuum Cumulant Expansion}
Let \(C_r\) denote the complete weight of one connected vacuum-graph type
\(r\), including its coupling, integrals, and inverse symmetry factor.
A vacuum diagram containing \(n_r\) indistinguishable copies of each type
has weight
\(\prod_r C_r^{n_r}/n_r!\).  Summing over all multiplicities gives
\[
 Z[0]=\prod_r\sum_{n_r=0}^{\infty}\frac{C_r^{n_r}}{n_r!}
 =\prod_r e^{C_r}
 =\exp\left(\sum_r C_r\right).
\]
Therefore
\[
 \log Z[0]=\sum_r C_r.
\]
The factorials are precisely those obtained by permuting identical
connected components.  Thus the logarithm retains each connected vacuum
graph once, with the symmetry factor already contained in \(C_r\).

\SupplementarySolution{18}{Cancellation of Arbitrary Vacuum Bubbles}
Every diagram contributing to the unnormalized insertion functional can be
decomposed uniquely into a part \(D\) containing all external insertions
and a multiset \(V\) of vacuum components.  The sum over \(V\) is independent
of the external points and equals the vacuum functional:
\[
 \widetilde G_n(x_1,\ldots,x_n)
 =Z[0]\,G_n^{\rm no\ vac}(x_1,\ldots,x_n).
\]
Dividing by \(Z[0]\) therefore gives
\[
 G_n=\frac{\widetilde G_n}{Z[0]}
 =G_n^{\rm no\ vac}.
\]
This argument uses only component factorization and the exponential
counting of indistinguishable vacuum graphs, not the form or order of the
polynomial interaction.  It consequently holds to every perturbative
order.

\SupplementarySolution{19}{A Wave Functional as a Euclidean Boundary}
Insert a field basis at the initial boundary:
\[
 \braket{\phi_f}{e^{-TH}\psi_i}
 =\int\mathcal D\phi_i\,K_E[\phi_f,T;\phi_i,0]\Psi_i[\phi_i]
 =\int_{\phi(T)=\phi_f}\mathcal D\phi\,
 e^{-S_E[\phi]}\Psi_i[\phi(0)].
\]
The spectral expansion is
\[
 \sum_n e^{-E_nT}\Psi_n[\phi_f]
 \braket{n}{\psi_i}.
\]
If the initial state has nonzero vacuum overlap, then as \(T\to\infty\)
\[
 e^{E_0T}\braket{\phi_f}{e^{-TH}\psi_i}
 \longrightarrow
 \Psi_0[\phi_f]\braket{\Omega}{\psi_i}.
\]
Thus a semi-infinite Euclidean path integral prepares the ground-state wave
functional, while changing its far-boundary weight changes only the
overlap constant.

\SupplementarySolution{20}{A Parity-Twisted Euclidean Trace}
Since \(P\ket x=\ket{-x}\),
\[
 \operatorname{Tr}(Pe^{-TH})
 =\int dx\,\braket{x}{Pe^{-TH}x}
 =\int dx\,\braket{-x}{e^{-TH}x}.
\]
The corresponding Euclidean paths obey the twisted boundary condition
\[
 x(T)=-x(0).
\]
For the oscillator, parity of the \(n\)th energy eigenstate is
\((-1)^n\), so
\[
 \operatorname{Tr}(Pe^{-TH})
 =\sum_{n=0}^{\infty}(-1)^n
 e^{-\omega T(n+1/2)}
 =\frac1{2\cosh(\omega T/2)}.
\]
Integrating the quoted oscillator heat kernel with endpoints \(x\) and
\(-x\) gives the same Gaussian result, verifying the twisted path-integral
boundary condition.

\SupplementarySolution{21}{Pfaffians under a Grassmann Change of Basis}
For \(\theta=N\theta'\), the top Grassmann monomial transforms as
\(\theta_1\cdots\theta_{2N}=\det(N)\theta'_1\cdots\theta'_{2N}\).
Berezin integration must therefore transform inversely:
\[
 d^{2N}\theta=\det(N)^{-1}d^{2N}\theta'.
\]
Now evaluate
\[
 \int d^{2N}\theta\,
 e^{\frac12\theta^{\mathsf T}A\theta}
 =\operatorname{Pf}(A)
\]
after the change of variables.  The exponent contains
\(N^{\mathsf T}AN\), while the measure contributes the inverse determinant,
so
\[
 \operatorname{Pf}(A)
 =\det(N)^{-1}\operatorname{Pf}(N^{\mathsf T}AN).
\]
Rearrangement proves
\[
 \boxed{\operatorname{Pf}(N^{\mathsf T}AN)
 =\det(N)\operatorname{Pf}(A)}.
\]

\SupplementarySolution{22}{An Exact Four-Grassmann Interaction}
Choose \(d^4\theta=d\theta_4d\theta_3d\theta_2d\theta_1\),
\(\int d^4\theta\,\theta_1\theta_2\theta_3\theta_4=1\), and write
\[
 S=\frac12A_{ab}\theta_a\theta_b
 +\lambda\theta_1\theta_2\theta_3\theta_4.
\]
Only the top monomial contributes.  Since the quadratic term squared
supplies \(\operatorname{Pf}(A)\theta_1\theta_2\theta_3\theta_4\),
\[
 Z=\int d^4\theta\,e^{-S}
 =\operatorname{Pf}(A)-\lambda.
\]
No higher term exists because every Grassmann variable squares to zero.
In diagram language, \(\operatorname{Pf}(A)\) is the sum over the three
signed Gaussian pairings, while the single quartic vacuum vertex has
weight \(-\lambda\).  There are no multi-vertex graphs, so the diagrammatic
enumeration is already exact.

\SupplementarySolution{23}{Bosonic and Fermionic Determinant Powers}
For a positive real symmetric \(K\),
\[
 \int d^Nx\,e^{-\frac12x^{\mathsf T}Kx}
 =(2\pi)^{N/2}(\det K)^{-1/2}.
\]
For independent complex Grassmann pairs and an arbitrary invertible \(A\),
\[
 \int d^N\bar\theta\,d^N\theta\,
 e^{-\bar\theta A\theta}=\det A.
\]
Finally, for \(2N\) real Grassmann variables and antisymmetric \(B\),
\[
 \int d^{2N}\theta\,
 e^{\frac12\theta^{\mathsf T}B\theta}
 =\operatorname{Pf}(B),\qquad
 \operatorname{Pf}(B)^2=\det B.
\]
Commuting Gaussians integrate inverse widths, whereas Berezin integration
extracts the highest polynomial coefficient; this reverses the determinant
power.  Antisymmetry of \(B\) is required because the symmetric part
multiplies an antisymmetric variable product and vanishes.

\SupplementarySolution{24}{Fermionic Wick's Theorem from a Pfaffian}
For antisymmetric invertible \(A\), introduce real Grassmann sources:
\[
 Z[\eta]=\int d^{2N}\theta\,
 e^{\frac12\theta^{\mathsf T}A\theta+\eta^{\mathsf T}\theta}
 =\operatorname{Pf}(A)
 e^{\frac12\eta^{\mathsf T}A^{-1}\eta},
\]
with the source and differential orders held fixed.  Four left
derivatives give
\[
 \langle\theta_a\theta_b\theta_c\theta_d\rangle
 =(A^{-1})_{ab}(A^{-1})_{cd}
 -(A^{-1})_{ac}(A^{-1})_{bd}
 +(A^{-1})_{ad}(A^{-1})_{bc}.
\]
The middle pairing is odd because forming it requires one interchange of
Grassmann objects.  With \(2N\) insertions, differentiation produces the
Pfaffian of the restricted covariance matrix, equivalently the signed sum
over all complete pairings.

\SupplementarySolution{25}{The One-Pair Berezin Gaussian}
Fix
\[
 \int d\bar\theta\,d\theta\,\theta\bar\theta=1.
\]
Because \((\bar\theta\theta)^2=0\),
\[
 e^{-\bar\theta a\theta}
 =1-\bar\theta a\theta
 =1+a\theta\bar\theta,
\]
and hence
\[
 \int d\bar\theta\,d\theta\,
 e^{-\bar\theta a\theta}=a.
\]
Under \(\theta\to\theta+\xi\) and
\(\bar\theta\to\bar\theta+\bar\xi\), all new terms have degree lower than
two in the integration variables or differ by a Berezin total derivative.
Their integrals vanish.  Thus the integral is translation invariant even
though Berezin integration is algebraic differentiation rather than an
ordinary limiting sum.

\SupplementarySolution{26}{The Complex Grassmann Determinant}
Change variables so that \(A\) is triangular; Berezin measures transform
with the inverse Jacobian, so the result is basis independent.  For a
diagonal kernel,
\[
 \prod_{i=1}^N\int d\bar\theta_i\,d\theta_i\,
 e^{-a_i\bar\theta_i\theta_i}
 =\prod_i a_i=\det A.
\]
Therefore, generally,
\[
 \boxed{\quad
 \int\prod_i(d\bar\theta_i\,d\theta_i)\,
 e^{-\bar\theta_iA_{ij}\theta_j}=\det A.
 \quad}
\]
The exponential terminates, and the integral selects the coefficient
containing every \(\theta_i\) and \(\bar\theta_i\).  That coefficient is the
antisymmetrized product defining \(\det A\), rather than its reciprocal.

\SupplementarySolution{27}{An Inverse Matrix from a Fermionic Insertion}
Differentiate the identity \(Z(A)=\det A\):
\[
 \frac{\partial Z}{\partial A_{lk}}
 =\det A\,(A^{-1})_{kl}.
\]
On the integral side,
\[
 \frac{\partial Z}{\partial A_{lk}}
 =\int d\bar\theta\,d\theta\,
 (-\bar\theta_l\theta_k)e^{-\bar\theta A\theta}
 =\int d\bar\theta\,d\theta\,
 \theta_k\bar\theta_l e^{-\bar\theta A\theta}.
\]
Consequently
\[
 \boxed{\quad
 \int d\bar\theta\,d\theta\,
 \theta_k\bar\theta_l e^{-\bar\theta A\theta}
 =\det A\,(A^{-1})_{kl}.
 \quad}
\]
The transpose-looking order is fixed by differentiating \(A_{lk}\), not
by a convention imposed afterward.

\SupplementarySolution{28}{The Four-Fermion Gaussian Sign}
Introduce sources and complete the square:
\[
 \frac{Z[\bar\eta,\eta]}{Z[0]}
 =\exp(\bar\eta A^{-1}\eta).
\]
Taking ordered source derivatives corresponding to
\(\theta_i\bar\theta_j\theta_k\bar\theta_l\) yields
\[
 \boxed{\quad
 \langle\theta_i\bar\theta_j\theta_k\bar\theta_l\rangle
 =(A^{-1})_{ij}(A^{-1})_{kl}
 -(A^{-1})_{il}(A^{-1})_{kj}.
 \quad}
\]
The second contraction crosses the middle pair.  Moving the associated
Grassmann derivatives next to each other requires one odd interchange,
which supplies the minus sign.  This is the finite-dimensional origin of
the sign rule for interchanging fermion lines.

\SupplementarySolution{29}{The Free Dirac Source Functional}
Let
\[
 \mathcal D(p)=i\gamma^\mu p_\mu+m-i0
\]
and use the source convention
\[
 -i\bar\psi\mathcal D\psi-i\bar\psi\eta-i\bar\eta\psi
\]
in the exponent.  The shifts
\(\psi\to\psi+\mathcal D^{-1}\eta\) and
\(\bar\psi\to\bar\psi+\bar\eta\mathcal D^{-1}\) give
\[
 \frac{Z[\bar\eta,\eta]}{Z[0]}
 =\exp\left(i\int\frac{d^4p}{(2\pi)^4}
 \bar\eta(-p)\mathcal D^{-1}(p)\eta(p)\right).
\]
Since
\[
 \mathcal D^{-1}(p)=
 \frac{-i\gamma\cdot p+m}
 {p^2+m^2-i0},
\]
ordered source differentiation reproduces
\[
 \langle T\psi(p)\bar\psi(-p)\rangle
 =-i\,\frac{-i\gamma\cdot p+m}
 {p^2+m^2-i0}
\]
for the displayed exponent convention.  The pole is fixed by the
fermionic vacuum wave functional just as in the scalar case.

\SupplementarySolution{30}{Fermion Determinants and Cyclic Signs}
For a background bilinear, the fermionic integral is
\[
 Z[\Gamma]\propto\det(\mathcal D+\Gamma)
 =\det\mathcal D\,
 \exp\left[\operatorname{Tr}\log
 (1+\mathcal D^{-1}\Gamma)\right].
\]
Expanding,
\[
 \operatorname{Tr}\log(1+\mathcal D^{-1}\Gamma)
 =\sum_{n=1}^{\infty}\frac{(-1)^{n+1}}n
 \operatorname{Tr}(\mathcal D^{-1}\Gamma)^n.
\]
The trace closes the chain of \(n\) propagators and vertices into one loop;
\(1/n\) removes cyclic choices of its starting vertex.  Relative to the
bosonic expansion, the alternating determinant series differs by one
overall minus sign for every connected closed chain.  Equivalently, closing
the chain requires an odd cyclic reordering of Grassmann fields.  Thus each
closed fermion loop carries the standard extra factor \(-1\).
\end{enumerate}
